%% GUS/MSS workshop paper draft.
%% Compile on Overleaf: pdfLaTeX + BibTeX (acmart is preinstalled).
%% Two-column workshop format (HotOS/PaPoC shape); adjust class options
%% to the venue CFP once chosen.
%%
%% Places that need the author's input are marked \todoinput{...}.

\begin{filecontents*}[overwrite]{gus.bib}
@article{gay2005subtyping,
  author  = {Gay, Simon and Hole, Malcolm},
  title   = {Subtyping for session types in the pi calculus},
  journal = {Acta Informatica},
  volume  = {42},
  number  = {2--3},
  pages   = {191--225},
  year    = {2005}
}
@inproceedings{honda2008multiparty,
  author    = {Honda, Kohei and Yoshida, Nobuko and Carbone, Marco},
  title     = {Multiparty Asynchronous Session Types},
  booktitle = {Proc.\ POPL},
  year      = {2008}
}
@article{dowling1984linear,
  author  = {Dowling, William F. and Gallier, Jean H.},
  title   = {Linear-time algorithms for testing the satisfiability of propositional {Horn} formulae},
  journal = {Journal of Logic Programming},
  volume  = {1},
  number  = {3},
  pages   = {267--284},
  year    = {1984}
}
@article{khanna2001approximability,
  author  = {Khanna, Sanjeev and Sudan, Madhu and Trevisan, Luca and Williamson, David P.},
  title   = {The Approximability of Constraint Satisfaction Problems},
  journal = {SIAM Journal on Computing},
  volume  = {30},
  number  = {6},
  pages   = {1863--1920},
  year    = {2001}
}
@inproceedings{ajmani2006modular,
  author    = {Ajmani, Sameer and Liskov, Barbara and Shrira, Liuba},
  title     = {Modular Software Upgrades for Distributed Systems},
  booktitle = {Proc.\ ECOOP},
  year      = {2006}
}
@inproceedings{zhang2021understanding,
  author    = {Zhang, Yongle and Yang, Junwen and Jin, Zhuqi and Sethi, Utsav and Rodrigues, Kirk and Lu, Shan and Yuan, Ding},
  title     = {Understanding and Detecting Software Upgrade Failures in Distributed Systems},
  booktitle = {Proc.\ SOSP},
  year      = {2021}
}
@inproceedings{alvaro2015lineage,
  author    = {Alvaro, Peter and Rosen, Joshua and Hellerstein, Joseph M.},
  title     = {Lineage-driven Fault Injection},
  booktitle = {Proc.\ SIGMOD},
  year      = {2015}
}
@article{litt2020cambria,
  author  = {Litt, Geoffrey and van Hardenberg, Peter and Good, Orion Henry},
  title   = {Project Cambria: Translate your data with lenses},
  journal = {Ink \& Switch technical report},
  year    = {2020},
  note    = {\url{https://www.inkandswitch.com/cambria/}}
}
@inproceedings{seco2020robust,
  author    = {Costa Seco, Jo{\~a}o and Ferreira, Paulo and Louren{\c c}o, Hugo and Ferreira, Carla and Ferr{\~a}o, Lu{\'i}s},
  title     = {Robust Contract Evolution in a TypeSafe MicroServices Architecture},
  booktitle = {The Art, Science, and Engineering of Programming},
  year      = {2020}
}
@inproceedings{raemaekers2014semantic,
  author    = {Raemaekers, Steven and van Deursen, Arie and Visser, Joost},
  title     = {Semantic Versioning versus Breaking Changes: A Study of the {Maven} Repository},
  booktitle = {Proc.\ SCAM},
  year      = {2014}
}
@article{lercher2024microservice,
  author  = {Lercher, Alexander and Glock, Johann and Macho, Christian and Pinzger, Martin},
  title   = {Microservice {API} Evolution in Practice: A Study on Strategies and Challenges},
  journal = {Journal of Systems and Software},
  volume  = {215},
  year    = {2024}
}
@misc{confluentcompat,
  author = {{Confluent, Inc.}},
  year   = {2026},
  title  = {Schema Evolution and Compatibility (Confluent Schema Registry)},
  note   = {\url{https://docs.confluent.io/platform/current/schema-registry/fundamentals/schema-evolution.html}}
}
@misc{pactcanideploy,
  author = {{Pact Foundation}},
  year   = {2026},
  title  = {Can {I} Deploy (Pact Broker documentation)},
  note   = {\url{https://docs.pact.io/pact_broker/can_i_deploy}}
}
@misc{bufbreaking,
  author = {{Buf Technologies}},
  year   = {2026},
  title  = {Buf breaking change detection rules},
  note   = {\url{https://buf.build/docs/breaking/rules/}}
}
@misc{oasdiff,
  author = {{Tufin (oasdiff maintainers)}},
  year   = {2026},
  title  = {oasdiff: {OpenAPI} diff and breaking changes},
  note   = {\url{https://github.com/oasdiff/oasdiff}}
}
@misc{serviceweaver,
  author = {{Google (Service Weaver team)}},
  year   = {2026},
  title  = {Service Weaver documentation: versioned rollouts},
  note   = {\url{https://serviceweaver.dev/docs.html}}
}
@misc{awsrollback,
  author = {Pokkunuri, Sandeep},
  year   = {2020},
  title  = {Ensuring rollback safety during deployments},
  note   = {Amazon Builders' Library. \url{https://aws.amazon.com/builders-library/ensuring-rollback-safety-during-deployments/}}
}
@misc{protobufrequired,
  author = {{Google (Protocol Buffers team)}},
  year   = {2026},
  title  = {Protocol Buffers Dos and Don'ts},
  note   = {\url{https://protobuf.dev/best-practices/dos-donts/}}
}
@misc{apicontextdrift,
  author = {{APIContext}},
  year   = {2026},
  title  = {The {API} Drift White Paper},
  note   = {\url{https://apicontext.com/resources/api-drift-white-paper/}}
}
\end{filecontents*}

\documentclass[sigconf,nonacm]{acmart}
\settopmatter{printacmref=false}

\usepackage{booktabs}
\usepackage{xcolor}
\usepackage{tcolorbox}

\newcommand{\todoinput}[1]{\begin{tcolorbox}[colback=orange!12,colframe=orange!70!black,title=\textbf{needs your input}]\small #1\end{tcolorbox}}

\newcommand{\Send}{\mathsf{Send}}
\newcommand{\Accept}{\mathsf{Accept}}
\newcommand{\Return}{\mathsf{Return}}
\newcommand{\Expect}{\mathsf{Expect}}
\newcommand{\leqreq}{\leq_{\mathsf{req}}}
\newcommand{\leqres}{\leq_{\mathsf{res}}}
\newcommand{\GUS}{\mathrm{GUS}}

\begin{document}

\title{What Can We Deploy Today?\\ Checking Upgrade Batches Against Every Version Pair}

\author{[AUTHOR NAME]}
\affiliation{\institution{[AFFILIATION]}\city{}\country{}}
\email{[EMAIL]}

\begin{abstract}
Schema-compatibility tools check one boundary between two versions.
Rolling deployments break that model in three ways: several services
upgrade as a batch, old and new instances of each service run
concurrently while the batch rolls, and some invariants span chains of
services rather than single call edges. We state the resulting problem
plainly: given a mesh of versioned, typed interfaces and a proposed
batch, is every version pair that can be live on any edge during the
rollout compatible? When the answer is no, the useful output is not a
verdict but a plan --- the subset of the batch that can ship, and the
order in which to ship it. Compatibility violations translate into
Horn clauses paired with ordering constraints. For the implication
fragment of that system a unique largest subset exists and unit
propagation finds it in linear time; conflicts --- target-state
failures and ordering cycles --- make maximality NP-hard and are
excluded conservatively. Every plan is
certified by replaying its stages. Finally, we show that batch
checking is still blind along the time axis: a guarantee can be
weakened in a rollout where nothing depends on it, and the violation
appears only when a dependent arrives several rollouts later, where a
per-rollout checker necessarily blames the wrong change. A small
ledger of guarantees, carried across rollouts, is enough to name the
rollout that weakened the guarantee. We develop the ideas
qualitatively on a nine-service mesh evolved through eleven synthetic
but plausible rollouts.
\end{abstract}

\maketitle

\todoinput{Author block above: name, affiliation, email. Also pick the
title --- alternatives: ``Ship What You Can: Checking Microservice
Upgrade Batches''; ``Upgrade Batches, Version Pairs, and Rollout
Order''. Delete the \texttt{nonacm} class option if the venue supplies
rights info.}

%% ---------------------------------------------------------------
\section{Introduction}
\label{sec:intro}

During a rolling upgrade a mesh runs a mixture of versions. Every
request may be served by an old or a new instance of its provider, and
issued by an old or a new instance of its caller. A change that is
compatible version-to-version can still break one of these transient
pairings, and the pairing occurs only in production. This failure
class is well documented: a study of 123 upgrade failures across eight
widely used distributed systems attributes about two thirds to
incompatible interaction between two coexisting versions, 40\% of
those through network messages~\cite{zhang2021understanding}.

Existing tools check pairs of versions at single boundaries. Pact
verifies consumer--provider contracts and gates deployment on recorded
verifications~\cite{pactcanideploy}; Buf classifies protobuf diffs by
wire compatibility~\cite{bufbreaking}; schema registries enforce
backward/forward compatibility per topic~\cite{confluentcompat};
\texttt{oasdiff} classifies OpenAPI diffs~\cite{oasdiff}. These are
sound answers to a smaller question. Three things are missing:

\begin{enumerate}
\item \textbf{Batches.} Per-boundary verdicts do not compose into a
  verdict for a set of simultaneous upgrades, and no per-boundary tool
  answers the operative question when a batch fails: which subset can
  ship?
\item \textbf{Order.} Many ``breaking'' changes are only
  \emph{order}-constrained. A provider that starts requiring a field
  its callers have agreed to send is safe exactly when every caller
  finishes rolling before the provider starts. The correct output is a
  schedule, and pairwise tools cannot express one.
\item \textbf{History.} A guarantee can be weakened in a rollout where
  nothing depends on it; every check passes and the change ships. When
  a dependent appears several rollouts later, the check that finally
  fails sits in the wrong diff. Attribution requires state that
  survives across rollouts.
\end{enumerate}

\begin{tcolorbox}[colback=black!4,colframe=black!60,title=\textbf{Problem statement}]
\small
Given a mesh of services with versioned, typed interfaces, a current
deployment state, and a proposed upgrade batch:
\textbf{(1)}~decide whether every version pair that can be live on any
edge during the rollout is compatible, and whether declared
cross-service dataflow chains hold in the resulting state;
\textbf{(2)}~if not, produce a subset of the batch and a stage order
that pass the same checks, certified by replaying the stages;
\textbf{(3)}~when a chain violation is exposed by a new dependent,
identify the earlier rollout that weakened the guarantee.
\end{tcolorbox}

We built a proof-of-concept checker, GUS (\emph{global upgrade
safety}), over OpenAPI interfaces, and evolved a nine-service mesh
through eleven synthetic rollouts to exercise it. The contribution of
this paper is the problem formulation and the shape of the solution:
compatibility as a check over all live version pairs
(\S\ref{sec:model}--\ref{sec:edge}), plans as a Horn system with
ordering constraints and a replay certificate (\S\ref{sec:mss}), and a
guarantee ledger for attribution across rollouts
(\S\ref{sec:ledger}). All evaluation is qualitative
(\S\ref{sec:casestudy}); \S\ref{sec:limits} states what the approach
cannot see.

\todoinput{One paragraph only you can write: who runs a reviewed batch
in practice (platform teams shipping umbrella releases, monorepo
deploy trains, on-prem vendors), and why that is your target user.
Teams that deploy each service independently and continuously have no
batch, so this framing matters for the paper's relevance claim.}

%% ---------------------------------------------------------------
\section{Model}
\label{sec:model}

A mesh is a directed graph $G = (S, E)$; an edge $e = (u \to v)$ is an
RPC endpoint on provider $v$ called by $u$. Each service has totally
ordered versions. A \emph{deployment state} is a map
$\theta : S \to \mathbb{V}$; an upgrade batch is a partial map $U$,
and the target state is $\theta' = \theta \oplus U$.

Each version of a service fixes four schemas per edge it participates
in: the caller's request body ($\Send$) and response expectation
($\Expect$), and the provider's request schema ($\Accept$) and
response body ($\Return$). Callers and providers are typed separately
because caller code drifts from provider contracts. When a caller
declares no outbound contract, the checker anchors
$\Send := \Accept(e,\theta)$ and $\Expect := \Return(e,\theta)$ at the
\emph{old} provider version --- the caller is assumed to have been
built against the contract that was live --- which reduces the check
to ordinary backward-compatibility rather than inventing caller-side
change.

Rolling semantics: replacing instances service by service means that,
at any instant, each upgrading service has live instances at
$\theta(s)$, at $\theta'(s)$, or both. We assume at most two versions
of a service are ever live (no canary with three), and that the mesh
routes a request to any live provider instance. Consequently the
version pairs that can be live on edge $u \to v$ during an
unconstrained rollout of $U$ are exactly
$\{\theta(u), \theta'(u)\} \times \{\theta(v), \theta'(v)\}$.
A \emph{staged} rollout restricts which pairs occur: if $u$ finishes
before $v$ starts, the pair $(\theta(u), \theta'(v))$ never exists.

%% ---------------------------------------------------------------
\section{Edge compatibility over version pairs}
\label{sec:edge}

Compatibility of one version pair is subtyping with the standard
variance split, as in binary session
subtyping~\cite{gay2005subtyping}. Writing $i$ for a caller version
and $j$ for a provider version on edge $e$:
\[
\mathrm{ok}_e(i,j) \;\triangleq\;
\Send_e(i) \leqreq \Accept_e(j)
\;\wedge\;
\Return_e(j) \leqres \Expect_e(i).
\]
The request relation $\leqreq$ requires everything the caller may send
to be admitted by the provider; the response relation $\leqres$
requires everything the provider may return to be understood by the
caller. Both are structural over a JSON type language: primitives
under a partial order, enums as value sets, records with
width and optionality rules, unions with existential matching, and
equi-recursive references handled by cycle detection with a
coinductive assumption at back-edges. The primitive order is strict:
an integer is \emph{not} assumed readable where a string is declared,
because production decoders reject such tokens. Findings that are
range risks rather than shape breaks (e.g.\ a response widening from
32- to 64-bit integers) are reported as warnings and do not fail a
batch.

Edge compatibility for the transition $\theta \to \theta'$ is the
conjunction over all live pairs:
\[
\mathrm{EdgeOK}(e, \theta, \theta') \;\triangleq\;
\bigwedge_{\substack{i \,\in\, \{\theta(u),\,\theta'(u)\}\\
                     j \,\in\, \{\theta(v),\,\theta'(v)\}}}
\mathrm{ok}_e(i,j).
\]
The $(\theta,\theta)$ conjunct is the precondition that the current
state is internally consistent, checked separately; the two mixed
conjuncts and the $(\theta',\theta')$ conjunct are what the transition
adds. $\GUS(\theta, U)$ is $\mathrm{EdgeOK}$ over all edges touched by
$U$, together with chain integrity in the target state
(\S\ref{sec:ledger}); a chain already broken in the baseline is
reported as pre-existing rather than attributed to the batch.

\begin{proposition}[Adequacy for edges]
\label{prop:adequacy}
Assume at most two live versions per service and uniform routing, and
that $\theta$ is consistent. The batch $U$, rolled without ordering
constraints, exposes no incompatible (caller, provider) version pair
on any edge iff $\mathrm{EdgeOK}(e,\theta,\theta')$ holds for every
edge $e$.
\end{proposition}
\begin{proof}[Proof sketch]
Each direction is immediate from the pair enumeration above: every
reachable pair is one of the four conjuncts, and each mixed or target
pair is realized by some interleaving (roll only $v$ partway; only
$u$; complete both).
\end{proof}

We scope the proposition deliberately: it speaks of edge-local pairs
only. Chain integrity (\S\ref{sec:ledger}) is checked in the endpoint
states $\theta$ and $\theta'$, not in transient mixtures; a chain that
breaks only mid-rollout is outside the current formulation, and we
flag this honestly in \S\ref{sec:limits}.

The value of checking mixed pairs is that the two relations are
asymmetric, so the same change can break one pair and leave the other
intact. In the case study, a provider adds a value to a response enum
while its caller upgrades in the same batch. The pair (old caller, new
provider) fails --- an old caller can receive the new value and has no
case for it --- while the pair (new caller, old provider) passes,
since old provider instances return only the original values. The violation is therefore not ``this change is
breaking'' but ``these two versions must not coexist,'' which is an
ordering fact; the next section turns such facts into plans.

%% ---------------------------------------------------------------
\section{From violations to plans}
\label{sec:mss}

When $\GUS(\theta,U)$ fails, the operative question is what subset of
$U$ can ship, and how. Define a \emph{plan} as a subset
$U' \subseteq U$ together with a partition of $U'$ into ordered stages.

\begin{definition}[Plan safety]
\label{def:plan}
A plan is safe from $\theta$ iff, taking $\theta_0 = \theta$ and
$\theta_k = \theta_{k-1} \oplus U'_k$ for stage $k$, each stage
satisfies $\GUS(\theta_{k-1}, U'_k)$.
\end{definition}

By Proposition~\ref{prop:adequacy} applied stage by stage, a safe plan
exposes no incompatible edge pair at any point of its execution, and
plan safety is directly checkable: \emph{replay} the stages, running
the full check against each accumulated intermediate state. The
checker always emits this certificate check; the solver below is a
candidate generator whose output is verified by replay, so its
soundness does not rest on the completeness of any encoding.

\paragraph{Generating candidates.}
Each violated conjunct pins one service of the batch at its new
version and one at its old version, and therefore reads as an ordering
fact. If the pair (old $u$, new $v$) is incompatible on either leg,
then new-$v$ instances must never coexist with old-$u$ instances: if
$u$'s contract cannot change ($u \notin U$, or $u$ has no declared
outbound contract), the violation is unconditional and $v$ is excluded
($\neg x_v$); otherwise $v$ may ship only if $u$ ships first --- the
definite clause $\neg x_v \vee x_u$ together with the precedence
$u \prec v$. Incompatibility of the symmetric pair (new $u$, old $v$)
swaps the roles. Incompatibility of the target pair
$(\theta',\theta')$ pins both services at their new versions; no order
avoids that pair, so the two upgrades conflict outright.

\begin{proposition}[Easy fragment]
\label{prop:easy}
For clause sets containing only definite clauses ($\neg a \vee b$) and
negative units ($\neg a$), the true-sets of satisfying assignments are
closed under union; hence a unique maximum exists, and propagating
falsity from the units computes it in $O(|\mathit{clauses}| + |U|)$.
\end{proposition}
\begin{proof}[Proof sketch]
If $a \in T_1 \cup T_2$ then $a \in T_i$ for some $i$, so any clause
$a \to b$ puts $b \in T_i \subseteq T_1 \cup T_2$; negative units hold
in both. The union of all models is then the maximum model.
Propagation is the classical linear-time procedure for Horn
theories~\cite{dowling1984linear}.
\end{proof}

\begin{proposition}[Hard boundary]
\label{prop:hard}
Once conflict clauses $\neg x_u \vee \neg x_v$ are admitted, deciding
a maximum-cardinality satisfying subset is NP-hard.
\end{proposition}
\begin{proof}[Proof sketch]
Such clauses are still Horn, so satisfiability stays linear; but
union-closure fails ($\{u\}$ and $\{v\}$ are models, $\{u,v\}$ is
not), and maximizing true variables with one conflict clause per graph
edge is maximum independent set. This is the Max-Ones classification
for this constraint family~\cite{khanna2001approximability}.
\end{proof}

The two propositions dictate the design. Unit propagation resolves the
easy fragment exactly. Conflicts --- target-pair failures, and cycles
in the precedence relation --- are resolved conservatively: the
checker excludes every member of a conflicting group and
re-propagates, rather than making an arbitrary choice or solving an
NP-hard optimization inside a safety tool. Cycles arise naturally. In
the case study, checkout starts requiring an idempotency key its old
callers do not send, so frontend must finish rolling first; in the
same version it changes a response field from string to integer while
frontend upgrades to expect the integer, so checkout must finish first
--- old checkout instances still answer with strings. The two
constraints are contradictory: no rolling order ships the pair, and
the checker excludes both. The result is sound but not maximum-cardinality, and the
tool says so. Surviving precedences are topologically sorted into the
plan's stages, and the plan is certified by replay
(Definition~\ref{def:plan}).

%% ---------------------------------------------------------------
\section{Chains and the guarantee ledger}
\label{sec:ledger}

Some invariants span more than one edge. In the case-study mesh,
checkout mints an order identifier that must arrive intact at the
email service; the intermediate schemas mark the field optional, which
is precisely why no edge-local check fires if checkout stops
guaranteeing it. The failure class is real: enforcing required fields
at intermediate hops of forwarding chains caused enough production
incidents that protobuf removed \texttt{required}
entirely~\cite{protobufrequired}.

The checker supports declared chains: a source annotates the field
that originates an identity (\texttt{x-provides}), a sink annotates
where it must arrive (\texttt{x-requires}), and intermediates annotate
only renames (\texttt{x-alias}). $\mathrm{chain\_ok}(K, \theta)$
enumerates the call paths from source to sink and requires the
identity to be present, non-null, and type-stable at the source and at
every intermediate hop, validated against what each hop \emph{sends
onward} (its outbound contract) --- which is what makes a rename
detectable at all. Because a mesh may route the identity along several
paths, every simple path between source and sink is checked.
Identity types must be stable end-to-end: an identity is stored and
joined on, not merely parsed.

\paragraph{Why one batch is not enough.}
Chain integrity at a state $\theta_k$ is a predicate on that state
alone, but the change responsible for a violation may sit in an
earlier rollout. The concrete sequence from the case study: rollout
$i$ makes checkout's \texttt{shipment\_ref} optional while no service
requires it --- there is no chain, nothing fails, the change ships;
rollout $k > i$ adds the requirement in email; the check at $k$ fails.
The only change a per-rollout checker can act on at step $k$ is
email's, and excluding email is what the solver correctly does for
that batch --- but the defect was introduced at step $i$, and the
repair belongs at checkout.

The fix is bookkeeping, not new theory. The checker maintains a
\emph{ledger}: after each rollout, for every provided identity ---
whether or not anything requires it yet --- it records the guarantee
as a tuple (provider service, field, type, required, nullable,
carrying paths), and appends an event when the tuple changes:
\textsf{born}, \textsf{eroded} (required became optional, or non-null
became nullable), \textsf{restored}, \textsf{demanded},
\textsf{violated}. The ledger is a JSON file persisted between
invocations; steps already recorded are skipped, so it accumulates
across rollouts exactly like the mesh it describes. When a chain check
fails, the ledger reports the most recent weakening event for that
identity alongside the failure. Attribution within the failing batch
is by reverting each on-path upgrade in isolation and re-checking; a
change whose revert repairs (or dissolves) the chain is marked a
culprit. This revert-and-recheck style of attribution is the static
analogue of the counterfactual reasoning used by lineage-driven fault
injection~\cite{alvaro2015lineage}, and we claim no more for it than
that.

For the sequence above the ledger output is (abridged):

{\small\begin{verbatim}
identity "shipment-tracking" -- EXPOSED
 born     @03: checkout provides shipment_ref
               (string, required)
 eroded   @07: required -> optional at checkout
 demanded @11: required by email -- REJECTED
 violated @11: chain-weakened
   -- guarantee last weakened at step 07
\end{verbatim}}

\noindent Step 11's batch is correctly blocked; step 07 is correctly
named. Neither statement is available to a checker that sees one
rollout at a time.

%% ---------------------------------------------------------------
\section{Case study}
\label{sec:casestudy}

We ported Google's Online Boutique demo to a nine-service mesh
(frontend, checkout, shipping, product catalog, currency, cart,
recommendation, payment, email; 14 RPC edges) with OpenAPI interfaces,
including caller-declared outbound contracts, and wrote an
\emph{evolution suite}: eleven rollout steps shaped like ordinary
feature work, where each step's baseline is what the previous step's
plan actually shipped. The scenarios are synthetic and constructed to
jointly cover the checker's rules; the suite doubles as its regression
oracle (expected verdicts, exact safe subsets, stage orders, and chain
outcomes are asserted). Table~\ref{tab:steps} summarizes; three steps
carry the argument.

\begin{table}
\caption{The eleven-step evolution suite. Bold rows are outcomes a
per-boundary, per-rollout checker cannot produce.}
\label{tab:steps}
\small
\begin{tabular}{@{}llp{3.6cm}@{}}
\toprule
Step & Outcome & Notes\\
\midrule
01 wishlist & ships & additive response fields\\
02 cart & blocked & response guarantee regressed\\
\textbf{03 accounts} & \textbf{ships in 2 stages} & callers first, then shipping\\
04 payments & blocked & response enum widened; nullable id\\
\textbf{05 payments} & \textbf{ships in 2 stages} & checkout before payment\\
06 money, expand & partial & widening ships as a warning; a map value change is blocked\\
\textbf{07 money, contract} & \textbf{ships in 2 stages} & four violation kinds, one order; quietly weakens a guarantee\\
\textbf{08 tracing} & \textbf{partial} & rename without alias: only the chain fails\\
09 tracing & ships & alias bridges the rename\\
\textbf{10 promo} & \textbf{blocked} & every mixed pair passes; only the target pair conflicts\\
\textbf{11 notifications} & \textbf{blocked} & step-07 weakening exposed; ledger names step 07\\
\bottomrule
\end{tabular}
\end{table}

\emph{Order (step 03).} Shipping begins requiring an
\texttt{account\_tier} field that both callers simultaneously upgrade
to send. Every pairwise tool classifies the provider change as
breaking. The checker instead emits a two-stage plan --- callers, then
shipping --- and the replay certificate passes: the only bad pair (old
caller, new provider) never exists under that order.

\emph{Target state (step 10).} Shipping accepts an optional
enum-valued promo code; frontend upgrades to always send a free-form
string. Both mixed pairs pass: old frontends omit the optional field,
and old shippings ignore the unknown field. Only the
$(\theta',\theta')$ pair conflicts. Registry-style
backward/forward checking relates successive versions of \emph{one}
schema; this failure lies between two parties' new schemas, a pair no
per-schema check ever relates --- here it is the entire failure.

\emph{History (steps 07 and 11).} As in \S\ref{sec:ledger}: the
weakening ships silently at step 07 because nothing requires the
identity; the requirement arrives at step 11 and is blocked; the
ledger names step 07. Without the ledger, step 11's failure is
attributed to email, which is the one service that did nothing wrong.

%% ---------------------------------------------------------------
\section{Related work}
\label{sec:related}

Per-boundary compatibility checking is mature. Pact's broker verifies
a candidate against the versions recorded as deployed --- a dynamic,
example-based check of mixed pairs for one consumer--provider
relationship~\cite{pactcanideploy}. Confluent's
\textsc{full\_transitive} mode is the both-directions condition for
pub/sub topics~\cite{confluentcompat}. Buf~\cite{bufbreaking} and
\texttt{oasdiff}~\cite{oasdiff} classify two-version diffs with
direction awareness. None of these composes verdicts across a mesh,
outputs an order, computes a shippable subset, or spans multiple hops.
Mixed-version operation has long been studied at the node level:
Ajmani et al.\ run simulation objects so nodes interoperate across
versions during asynchronous upgrades~\cite{ajmani2006modular};
Service Weaver avoids cross-version communication entirely by routing
each request within one application version~\cite{serviceweaver};
Amazon's rollback-safety guidance prescribes two-phase serialization
changes as an operational discipline~\cite{awsrollback}. Typed
evolution has been addressed with language support (contract evolution
with runtime adaptation~\cite{seco2020robust}) and with bidirectional
lenses that translate between schema versions~\cite{litt2020cambria}
--- repair, where this work does detection. The per-pair subtype
relations are those of binary session subtyping~\cite{gay2005subtyping};
we do not use multiparty session types~\cite{honda2008multiparty},
since no global protocol is checked here. The solver core is
Dowling--Gallier propagation~\cite{dowling1984linear}, and the
hardness boundary is the Max-Ones
classification~\cite{khanna2001approximability}. That declared
versions understate breakage is well
evidenced~\cite{raemaekers2014semantic,lercher2024microservice}.

%% ---------------------------------------------------------------
\section{Limitations}
\label{sec:limits}

\emph{Specs, not behavior.} The checker judges declared contracts.
Measured drift between production APIs and their OpenAPI documents is
substantial~\cite{apicontextdrift}, so a passing verdict means ``no
statically visible hazard in the declared mesh''; failing verdicts are
near-certain wire breaks. The asymmetry should be stated wherever the
tool is used. \emph{Chains in transient states.} Chain integrity is
checked in the endpoint states of each rollout, not in mid-rollout
mixtures (Proposition~\ref{prop:adequacy} deliberately excludes it);
extending the pair-enumeration argument to chains is open.
\emph{Two live versions.} Canary deployments hold three or more
versions; the pair enumeration generalizes, the tool does not yet.
\emph{Annotations.} Chains require hand-written
provides/requires/alias annotations and caller-declared outbound
contracts; stale annotations corrupt verdicts silently. Inferring both
from traffic or source is the obvious next layer of this work.
\emph{Conservative choices.} Conflict groups and culprit sets are
excluded wholesale where choosing among them is NP-hard; a weighted
formulation would let operators express preferences.
\emph{Response-carried identities.} Chain discovery follows forward
call edges, so identities carried in responses are out of scope.
\emph{Rollback.} An aborted rollout traverses the same pairs in
reverse; the model covers it by swapping $\theta$ and $\theta'$, the
tool does not yet.

\todoinput{Evaluation plan: commit to what you will actually do next
(replaying public OpenAPI version histories; reproducing documented
upgrade failures from the SOSP'21 taxonomy as mesh scenarios; or a
case study with a platform team's umbrella releases), and what you
will measure.}

\todoinput{Acknowledgments, venue footnote, and whether to link the
artifact repository.}

\bibliographystyle{ACM-Reference-Format}
\bibliography{gus}

\end{document}
